	\documentclass[12pt]{article}

		\usepackage{graphicx}
		\usepackage{amsmath,amssymb,rotating,multirow,setspace,fullpage,verbatim,subfigure,quotes,url,bm}

	\usepackage[left=1in,right=1in,top=1in,bottom=1.1in]{geometry}
		\usepackage[pdftex]{color}
		\makeatletter
		\g@addto@macro{\UrlBreaks}{\UrlOrds}
		\makeatother
		\usepackage{relsize}
		%\usepackage{hyperref}
		\usepackage{placeins}
		%\usepackage{csquotes}
		\usepackage{chronology}
		\usepackage{verbatim}
		\bibliographystyle{apsr}
		\clubpenalty=10000 \widowpenalty=10000
		\newcommand\vm[1]{\bm{\mathrm{#1}}}
		\newenvironment{proof}[1][Proof]{\noindent\textbf{#1.} }{\ \rule{0.5em}{0.5em}}
		\newtheorem{result}{Result}
		\newtheorem{lemma}{Lemma}
		\newtheorem{proposition}{Proposition}
		\newtheorem{corollary}{Corollary}
		\usepackage{fancyhdr}
		\usepackage{lscape}

		%\newcommand{\mmcomment}[1]{{\textcolor{blue}{\textsc{\textbf{[#1 --mm]}}}}}
		
		\newtheorem{listitem}{Listitem}
		% below from Elisha 
		\newtheorem{assumption}{Assumption}
		\newtheorem{definition}{Definiton}
		\newtheorem{theorem}{Theorem}
		\newtheorem{hypothesis}{Hypothesis}
		
		
		% from internet: https://tex.stackexchange.com/questions/100880/hypotheses-and-subhypotheses-with-ntheorem
		\usepackage{ntheorem}
		\newtheorem{hyp}{Hypothesis} 
		\makeatletter
		\newcounter{subhyp} 
		\let\savedc@hyp\c@hyp
		\newenvironment{subhyp}
		{%
			\setcounter{subhyp}{0}%
			\stepcounter{hyp}%
			\edef\saved@hyp{\thehyp}% Save the current value of hyp
			\let\c@hyp\c@subhyp     % Now hyp is subhyp
			\renewcommand{\thehyp}{\saved@hyp\alph{hyp}}%
		}
		{}
		\newcommand{\normhyp}{%
			\let\c@hyp\savedc@hyp % revert to the old one
			\renewcommand\thehyp{\arabic{hyp}}%
		} 
		\makeatother
		
		\thispagestyle{empty}

		%\fancyhf{}
	%	\rhead{}	
	%	\lhead{}

		\begin{document}

			\thispagestyle{fancy}
			\pagenumbering{arabic} 
				
		%	\tableofcontents{\fancyhead[LE]{\slshape\contentsname}}
			
		%	\newpage 
			
				\begin{center}
				{\large \textbf{QTM 220: Regression Analysis } }\\
				\vspace{0.1in}
				Emory University\\
				\vspace{0.2in}
				\begin{tabular}{llll}
				&	Lecture room:   & ***  &\\
					& Lecture time:   & *** & \\
					& Lab room: & & \\
				& 	Lab time: & & \\
					\\
					Instructors:     & Adam Glynn  & & Allison Cuttner  \\
					Email: & adam.glynn@emory.edu & &  allison.cuttner@emory.edu \\
					Office: & PAIS & & PAIS \\
					Office hours: &  & & \\
					\\
					& TA: & & \\
				& 	Email: & & \\
					
				\end{tabular}
			\end{center}
			
			
			\subsection*{Course Description}
(from previous syllabus) This course introduces students to widely-used procedures for regression analysis and provides intuitive, applied, and formal foundations for more advanced methods treated in later courses. The course covers descriptive, population and causal inference with regression from a modern perspective. While the course will emphasize the mathematical foundations of these concepts, each topic will also cover the implementation of the relevant methods in an application through the use of the statistical program R.

This course represents the culmination of the required courses for QTM majors. It uses the coding skills learned in QTM 150 and the math skill learned in Math 210 and QTM 220 to conduct much of the analysis discussed in QTM 110. QTM 220 also provides an introduction and a foundation for future study  of topics such as machine learning and causal inference. 


{\bf new things to thing about}: How does this course fit in to the curriculum/sequence for the rest of the major? Why is it important for students to devote a lot of time/effort to? Why are description, population inference, and causal inference separated? Why are different types of covariates separated? How will this course prepare students to do `cool' stuff like causal inference and machine learning? 
	
\bigskip


\paragraph{Course Goals} By the end of this course, students will be able to 

{\bf These are copied from old syllabus and need to be updated }

\begin{itemize}
\item Be able to conduct a number of regression analyses and correctly interpret the results based on research design concepts learned in QTM 110.
\item Understand the properties of regression estimators so as to understand when these estimators are reliable.
\item Distinguish between descriptive, population, and causal interpretations of regression findings, and understand the assumptions that are necessary for each.
\end{itemize}

{\bf new ideas:} \\
student will be able to...
\begin{itemize}
\item formalize questions about the world (descriptive, population, and causal)
\item use regression to answer those questions 
\item quantify uncertainty and error 
\item understand when certain techniques are appropriate to use (or not) 
\end{itemize}

What other tools/skills (including specific R skills) will we be confident that students will have mastered by the end of the course? 

What tools/skills/concepts will students be exposed to (but not master) by the end of the course? 

\bigskip

\paragraph{Prerequisites}

QTM 210 or equivalent is prerequisite. 

			\subsubsection*{Course Requirements} 
			\paragraph{Attendance} Lecture attendance is expected. If a student cannot attend lecture, they must email the instructor to explain their absence. Please do not come to class (lecture or lab) if sick. 
			
			\paragraph{Problem Sets} Problem sets will be assigned (approximately) biweekly throughout the semester. The homework assignments will consist of analytical problems, computer work, and/or data analysis. The homework write-up must be word processed in some manner (exact manner to be determined by the TAs). The due date and time for all homework will be indicated on each assignment and should be submitted through Canvas. No late work will be accepted without an approved excuse.

We encourage students to confer with each other on the homework assignments, but you must write your own solutions (this includes computer work). Copying and pasting code from other students is not allowed, and you must be able to explain anything you turn in.		
			
			\paragraph{Midterm Exams} The midterms will be in-class exams {\it  There is no collaboration allowed on the midterms.}

			\paragraph{Final Exam}  The final will be held during the assigned exam time and location. {\it  There is no collaboration allowed on the final.}  
			
			\paragraph{Lab Session} Lab sessions will be run by the TAs. The purpose of the lab sessions is to learn the R programming and data analysis skills useful for regression and other empirical analysis. Students must bring a laptop, with R installed and updated on it, to these sessions. Lab sessions will prepare students to complete the problem sets and review the solutions of previous homeworks. Lab participation and course performance are highly correlated. 
			
			\paragraph*{Policies} any other course policies??? laptops in lecture? using Canvas discussion threads to ask questions about course material instead of emailing? extra credit (and lack thereof)? grade appeals? Pablo's grade appeal policy: If you believe that your grade on any assignment or exam question is incorrect or unfair, you should submit your concerns, in writing, to the instructor.  There are two types of appeals: 
1.	Addition errors or other administrative snafus: Please let me know if you catch one of these ASAP.  I’ll confirm the error and change your grade in the gradebook. 
2.	Disagreement about grading criteria:  Please wait two days (or the nearest non-weekend day) after receiving your grade before submitting a written appeal.  Do not leave comments on Canvas assignments for TAs as they are not responsible for grade appeals or score corrections.  The written appeal should fully summarize what you believe the problems are and why.   The instructor will consider your appeal. If you are not satisfied with the response, you may resubmit the assignment for regrading in its entirety by the instructor  This grade will be final.  Note that grades may go up or down during an appeal.  Failure to comply with this appeal process means that you forfeit your appeal. 

		
	\subsubsection*{Grading }		
	
			\begin{itemize}
			\item Problem Sets (35\%)
			\item Midterm Exams (40\%)
			\item Final Exam (25\%)
			\end{itemize}
			
			\subsubsection*{Incomplete Grades}
			No incomplete grades will be given unless there is an agreement between the instructor and the student {\bf prior} to the end of the course.  The instructor retains the right to determine legitimate reasons for an incomplete grade.
			
\subsubsection*{Honor Code}

All students enrolled at Emory are expected to abide by the Emory College Honor Code. Any type of academic misconduct is not allowed which includes 1) receiving or giving information about the content or conduct of an examination knowing that the release of such information is not allowed and 2) plagiarizing, whether intentionally or unintentionally, in any assignment. For the activities that are considered to be academically dishonest, refer to the Honor Code: http://catalog.college.emory.edu/academic/policies-regulations/honor-code.html.\\

\subsubsection*{Accessibility and Accommodations}

If you are seeking classroom accommodations or academic adjustments under the Americans with Disabilities Act, you are required to register with Office of Accessibility Services (http://accessibility.emory.edu). To receive academic accommodations for this class, please obtain the relevant letter and meet with the instructor at the beginning of the semester. Students are expected to give two weeks notice of the need for accommodations.
%Emory University is committed under the Americans with Disabilities Act and its Amendments and Section 504 of the Rehabilitation Act to providing appropriate accommodations to individuals with documented disabilities. If you have a disability-related need for reasonable academic adjustments in this course, provide the instructor(s) with an accommodation notification letter from Access, Disabilities Services and Resources office (\href{http://www.ods.emory.edu/}{http://www.ods.emory.edu/}). Students are expected to give two weeks notice of the need for accommodations. If you need immediate accommodations or physical access, please arrange to meet with instructor as soon as your accommodations have been finalized.\\

\bigskip
			\newpage 
			\section*{Schedule}
\subsection*{Part 1}
Lecture 1 (Aug 23 | Aug 24)\\
Descriptive stats w 1 binary covariate Start w data, then versions of the question\\
Does everybody who goes to college make more (Lebron etc)\\
Do more people who go to college make at least x (fraction; number then cdf plot)\\
Do more people who go to college make (almost) exactly x (histogram)\\
Summarizing the histogram w/ one number: center\\
Median (people above = below)\\
Mean (money above = below)\\
A second number: spread\\
MAD\\
Variance\\

Means not as restrictive as you’d think — binary case is a frequency; 1 variable to frequency above threshold. Variance stuff can come last with computations for the binary case; plots. Use all CPS data.\\
Syllabus\\
Closing:\\
Nested means vs nested medians

\ \ 

Lecture 2 (Aug 28 | Aug 29)\\
Random sampling and sampling distributions within each value of the covariate separately. Why need? Efficiency (survey vs census, QC)\\
Why sensible?\\
Samples tend to agree with big thing.\\
(mean, histogram, etc)\\
Not exactly—sampling distn of the mean (mean, histogram, variance, concentration via Chebyshev)

\ \ 

Infinite Population Approx ?\\
(Random sub samples of cps look like random subsamples of random subsamples of cps)\\
Calculations with densities\\
Why? Helps us think, more or less

\ \ 

Lecture 3 (Aug 30 | Aug 31)\\
Sampling distns and confidence intervals within each value of the covariate and difference across the covariate.
(Chebyshev)\\
Coverage prob (conservative)

\ \ 

Universality\\
Exact via normal approximation, normality based\\
Repeat for differences

\ \ 

Lecture 4 (Sept 6 | Sept 5) order flips here: Adam + Allison now ahead of Skip\\
Causality and Prediction\\
Using regression to predict/impute missing outcome data Causal inference

\ \

Lecture 5 (Sept 11 | Sept 7)\\
1 Discrete Covariate (beyond binary)—description\\
Zoom in and it’s binary; new challenges are complex comparisons and sparsity\\
Solan to former: be clear. Examples!\\
Soln to the latter:\\
Specific: Grouping.\\
Problem: Misspecification\\
Finite: calculate effect of misspecification\\
Infinite pop approx: calculate effect of misspecification\\
Same for comparisons\\
Benign case

\ \ 

Lecture 6 (Sept 13 | Sept 12)\\
Misspecification and inference\\
Sampling dist w/ misspec\\
implications of misspec for CI

\ \

Lecture 7 (Sept 18 | Sept 14)\\
Causality\\
Binary comparision\\
Incremental effects\\
Treatment shift generally

\ \ 

Review (Sept 20 | Sept 19)

\ \ 

Exam (Sept 25 | Sept 21)



\subsection*{Part 2}
Lecture 8 (Sept 27 | Sept 26)\\
Beyond Bins 1 single ordered covariate \\
Least squares interpretation \\
Lines\\
Residual properties and plots\\
Variance calc\\
Sampling Distn

\ \ 

Lecture 9 (Oct 2 | Sept 28)\\
beyond bins 2 single ordered covariate\\
Quadratics, etc. Seasonality. Linalg. basis expansions--polynomials (taylor series) \\
Residual plots (uncorrelated interp)

\ \ 

Lecture 10 (Oct 4 | Oct 3)\\
single ordered covariate, population inference (means) \\
Cond means and least squares\\
Misspec and Uncorrelated interp for noise \\
(dist of linear functions of independent means)

\ \ 

Fall Break (Oct 9 | Oct 10)

\ \ 

Lecture 11 (Oct 11 | Oct 5)\\
single ordered covariate, pop inference (means) \\
Inference: distributions of summaries\\
Delta method\\
Implications of misspecification\\
IPW lsq

\ \ 

Lecture 12 (Oct 16 | Oct 12)\\
Continuum idea\\
Calculations in continuum

\ \ 

Review (Oct 18 | Oct 17)

\ \ 

Exam (Oct 23 | Oct 19)


\subsection*{Part 3}
Descriptive\\
Lecture 13 (Oct 25 | Oct 24) \\
multivariate discrete. like before but interactions \\
Saturated + OVB\\
Missing interaction bias

\ \ 

Lecture 14 (Oct 30 | Oct 26)\\
Multivar (X 1D continuous + D discrete)

\ \ 

Lecture 15 (Nov 1 | Oct 31)\\
Multivar (X 1D continuous + D continuous)

\ \ 

Lecture 16 (Nov 6 | Nov 2)\\
Multivar (X arbitrary + D binary/continous)\\
Hard to draw\\
Taylor series, additive, piecewise-constant\\
ML idea: honest trees

\ \ 

Election Day (Nov 8 | Nov 7)\\
It's really the 7th

\ \ 

Pop + Causal versions\\
Lecture 17 (Nov 13 | Nov 9 )\\
multivariate discrete. like before but interactions\\
Saturated + OVB \\
Missing interaction bias\\
Implications for summaries\\
IPW


\ \ 

Lecture 18 (Nov 15 | Nov 14)\\
Multivar (X 1D continuous + D discrete)

\ \ 

Lecture 19. (Nov 20 | Nov 16)\\
Multivar (X 1D continuous + D continuous)

\ \ 

Thanksgiving Break (Nov 22 | Nov 23)

\ \ 

Lecture 20 (Nov 27 | Nov 21)\\
Multivar (X arbitrary + D binary/continous)\\
Hard to draw

\ \ 

Lecture 21 (Nov 29 | Nov 28)\\
Pop version, multivar discrete

\ \ 

Lecture 22 (Dec 4 | Nov 30)\\
Pop version, X 1D continuous + D discrete

\ \ 

Review (Off-Schedule | Dec 5)

\ \ 

Final Exam (exam period) 		
		\end{document} 